\begin{table}[htbp]
\centering
\caption{Artigos sobre transferências intergovernamentais e sobre efeito \textit{flypaper} ou ilusão fiscal, por periódico: 2000--2025}
\label{tab:baker}
\small
\begin{tabular}{lccc}
\toprule
 & (1) & (2) & (3)\\
 & Transferências & \textit{Flypaper}/ilusão fiscal & (2)/(1) (\%)\\
\midrule
\textit{Public Finance Review} & 26 & 13 & 50,0\%\\
\textit{International Tax and Public Finance} & 24 & 5 & 20,8\%\\
\textit{Public Choice} & 22 & 14 & 63,6\%\\
\textit{Journal of Public Economics} & 12 & 3 & 25,0\%\\
\textit{National Tax Journal} & 7 & 4 & 57,1\%\\
\textbf{Finanças públicas e escolha pública} & \textbf{91} & \textbf{39} & \textbf{42,9\%}\\
\midrule
\textit{Public Budgeting \& Finance} & 17 & 10 & 58,8\%\\
\textit{Publius: The Journal of Federalism} & 15 & 3 & 20,0\%\\
\textit{Local Government Studies} & 12 & 7 & 58,3\%\\
\textit{Public Finance and Management} & 13 & 4 & 30,8\%\\
\textit{State and Local Government Review} & 7 & 1 & 14,3\%\\
\textbf{Federalismo, orçamento e governo local} & \textbf{64} & \textbf{25} & \textbf{39,1\%}\\
\midrule
\textit{Economia Aplicada} & 7 & 3 & 42,9\%\\
\textit{Revista de Administração Pública} & 6 & 0 & 0,0\%\\
\textit{Revista Econômica do Nordeste} & 6 & 1 & 16,7\%\\
\textit{Administração Pública e Gestão Social} & 5 & 4 & 80,0\%\\
\textit{Estudos Econômicos} & 4 & 2 & 50,0\%\\
\textbf{Periódicos brasileiros} & \textbf{28} & \textbf{10} & \textbf{35,7\%}\\
\midrule
\textbf{Total dos 15 periódicos} & \textbf{183} & \textbf{74} & \textbf{40,4\%}\\
Corpus completo (todos os periódicos) & 1.712 & 640 & 37,4\%\\
\bottomrule
\end{tabular}
\par\smallskip
\begin{minipage}{0.97\textwidth}\scriptsize
Nota: a Tabela~\ref{tab:baker} segue o formato da Tabela 1 de \citeonline{baker2022}. Resume o número de artigos publicados entre 2000 e 2025 em cinco periódicos de finanças públicas e escolha pública, cinco periódicos de federalismo, orçamento e governo local (os mais produtivos do corpus em cada grupo) e cinco periódicos editados no Brasil. O corpus foi obtido no OpenAlex com busca em título e resumo pelos termos ``intergovernmental transfers'', ``intergovernmental grants'', ``flypaper effect'' e ``fiscal illusion'' e pelos equivalentes em português e espanhol (``transferências intergovernamentais'', ``transferências constitucionais'', ``efeito flypaper'', ``ilusão fiscal'', ``transferencias intergubernamentales'', ``ilusión fiscal''), mantendo apenas artigos com autoria identificada e fonte primária classificada como periódico. A coluna (1) informa o total de artigos do corpus em cada periódico; a coluna (2), quantos deles tratam do efeito \textit{flypaper} ou da ilusão fiscal (termo presente no título, no resumo ou nas palavras-chave); a coluna (3), a participação de (2) em (1). Contagens produzidas com o pacote \textit{bibliometrix} \cite{aria2017} (\texttt{06\_codigos\_e\_planilhas/05\_bibliometria.R}).
\end{minipage}
\end{table}
